% Worked Examples: Concept 2.4.5 - Output Feedback Basics
\documentclass[11pt,a4paper]{article}
\usepackage{worked-examples-template}

\title{\textbf{Worked Examples}\\[0.5em]
\large Concept 2.4.5: Output Feedback Basics\\[0.3em]
\normalsize \coursesubtitle{} -- \coursetitle}
\author{Assoc.\ Prof.\ William Robertson \and Dr Sean McGowan}
\date{\today}

\begin{document}

\maketitle
\tableofcontents
\newpage

\section{Concept 2.4.5: Observer Design}

\subsection{Example 1: Full-State Observer Design}

\paragraph{Problem:} Consider a system where we can only measure the output, not all states:
\begin{align}
\frac{d}{dt}\begin{bmatrix} x_1 \\ x_2 \end{bmatrix} = \begin{bmatrix} 0 & 1 \\ -3 & -4 \end{bmatrix}\begin{bmatrix} x_1 \\ x_2 \end{bmatrix} + \begin{bmatrix} 0 \\ 2 \end{bmatrix}u, \quad y = \begin{bmatrix} 1 & 0 \end{bmatrix}\begin{bmatrix} x_1 \\ x_2 \end{bmatrix}
\end{align}

Design a full-state observer (Luenberger observer) with observer poles at $-5$ and $-6$ (faster than the plant dynamics).

\paragraph{Solution:}

The observer estimates the state $\hat{x}$ using the dynamics:
\begin{align}
\frac{d\hat{x}}{dt} = A\hat{x} + Bu + L(y - \hat{y})
\end{align}

where $L$ is the observer gain matrix and $\hat{y} = C\hat{x}$.

The error dynamics are $e = x - \hat{x}$:
\begin{align}
\frac{de}{dt} = (A - LC)e
\end{align}

We need to design $L$ so that $A - LC$ has eigenvalues at $-5$ and $-6$.

With $L = \begin{bmatrix} l_1 \\ l_2 \end{bmatrix}$ and $C = \begin{bmatrix} 1 & 0 \end{bmatrix}$:
\begin{align}
LC = \begin{bmatrix} l_1 \\ l_2 \end{bmatrix}\begin{bmatrix} 1 & 0 \end{bmatrix} = \begin{bmatrix} l_1 & 0 \\ l_2 & 0 \end{bmatrix}
\end{align}

Therefore:
\begin{align}
A - LC = \begin{bmatrix} 0 & 1 \\ -3 & -4 \end{bmatrix} - \begin{bmatrix} l_1 & 0 \\ l_2 & 0 \end{bmatrix} = \begin{bmatrix} -l_1 & 1 \\ -3-l_2 & -4 \end{bmatrix}
\end{align}

The characteristic polynomial of the observer error dynamics is:
\begin{align}
\det(sI - (A-LC)) &= \det\begin{bmatrix} s+l_1 & -1 \\ 3+l_2 & s+4 \end{bmatrix} \\
&= (s+l_1)(s+4) + (3+l_2) \\
&= s^2 + (4+l_1)s + 4l_1 + 3 + l_2
\end{align}

The desired characteristic polynomial with poles at $-5$ and $-6$ is:
\begin{align}
(s+5)(s+6) = s^2 + 11s + 30
\end{align}

Matching coefficients:
\begin{align}
4 + l_1 &= 11 \quad \Rightarrow \quad l_1 = 7 \\
4l_1 + 3 + l_2 &= 30 \\
28 + 3 + l_2 &= 30 \quad \Rightarrow \quad l_2 = -1
\end{align}

The observer gain is:
\begin{align}
L = \begin{bmatrix} 7 \\ -1 \end{bmatrix}
\end{align}

The complete observer equations are:
\begin{align}
\frac{d\hat{x}}{dt} = \begin{bmatrix} 0 & 1 \\ -3 & -4 \end{bmatrix}\hat{x} + \begin{bmatrix} 0 \\ 2 \end{bmatrix}u + \begin{bmatrix} 7 \\ -1 \end{bmatrix}(y - \hat{x}_1)
\end{align}

Expanding:
\begin{align}
\frac{d\hat{x}_1}{dt} &= \hat{x}_2 + 7(y - \hat{x}_1) \\
\frac{d\hat{x}_2}{dt} &= -3\hat{x}_1 - 4\hat{x}_2 + 2u - (y - \hat{x}_1)
\end{align}

\textbf{Interpretation:} The observer uses the output measurement $y$ to correct its state estimates. The observer poles at $-5$ and $-6$ ensure that estimation errors decay faster than the plant dynamics, providing rapid convergence of $\hat{x}$ to $x$.
\end{document}
